\section{Introduction}
ADMET (absorption, distribution, metabolism, excretion, toxicity) prediction is where much of
early drug discovery succeeds or fails, and Therapeutics Data Commons (TDC)
\citep{huang2021tdc,huang2022artificial} is the community-standard benchmark for it. Its
leaderboard culture optimizes a single metric under one fixed scaffold split---useful for ranking,
but answering a different question than the one a practitioner faces: \emph{if I deploy this model
on genuinely novel chemistry, how well will it work, and will it tell me when it does not know?}

A long line of work shows random splits leak structurally similar molecules across folds and
overstate prospective performance \citep{sheridan2013timesplit,wu2018moleculenet,deng2023systematic}.
Scaffold splits were introduced as a ``harder, more realistic'' default \citep{wu2018moleculenet}.
Yet recent results complicate this: \citet{fooladi2025evaluating} find that for ADMET, Bemis--Murcko
scaffold splits are \emph{not substantially harder than random}, and only cluster/distance splits
expose a real gap; \citet{guo2024scaffold} show scaffold splits still leak because different
scaffolds can be near-identical molecules. Separately, the uncertainty-quantification literature
offers conformal prediction \citep{tibshirani2019conformal,angelopoulos2023gentle}, deep ensembles
\citep{scalia2020evaluating,hirschfeld2020uncertainty}, and conformal-under-shift methods
\citep{laghuvarapu2023codrug,li2024conformalized}---but these are rarely tied back to the
split-induced optimism they could mitigate.

We are explicit that \emph{neither phenomenon is a new discovery}. Our contribution is a rigorous
\emph{consolidation and extension} that no single prior paper provides, and---in the spirit of
\citet{kapoor2023leakage}---one that reports negatives plainly. This extends the
trustworthy-ML-under-shift agenda of the Zitnik lab's own tools (TDC $\rightarrow$ TxGNN
\citep{huang2024txgnn}, the last of which defers uncertainty quantification to future work).

\paragraph{Contributions.}
(1) A statistically honest, per-endpoint \textbf{generalization-gap atlas} across 22 endpoints and
model families, engaging the scaffold$\approx$random result head-on and promoting cluster splits to
the primary stressor. (2) A mechanistic \textbf{applicability-domain} account of the gap. (3) An
\textbf{honest evaluation of uncertainty under shift}, including clean negative results, plus a
foundation-model arm asking whether pretrained representations close the gap.

\section{Related Work}
\textbf{Split realism.} \citet{sheridan2013timesplit} established that train--test dissimilarity
governs apparent accuracy; \citet{wu2018moleculenet} popularized scaffold splits;
\citet{deng2023systematic} confirmed scaffold\,$<$\,random performance. \citet{guo2024scaffold} and
\citet{fooladi2025evaluating} argue scaffold splits under-stress and cluster/distance splits are
needed; \citet{vantilborg2022exposing} expose a complementary within-scaffold failure (activity
cliffs). \textbf{TDC and critique.} TDC \citep{huang2021tdc,huang2022artificial} is our object of
study; concurrent work \citep{koleiev2026critical} audits the 22 ADMET leaderboards for leakage and
reproducibility---complementary to our controlled signed-gap measurement.
\textbf{Uncertainty / selective prediction.} We \emph{apply} (not invent) split conformal
\citep{tibshirani2019conformal,angelopoulos2023gentle}, deep ensembles
\citep{scalia2020evaluating,hirschfeld2020uncertainty}, and conformal-under-covariate-shift
\citep{laghuvarapu2023codrug,li2024conformalized}, evaluated via risk--coverage curves
\citep{geifman2017selective}. \textbf{Foundation models.} We evaluate frozen ChemBERTa
\citep{ahmad2022chemberta2} (cf.\ MolFormer \citep{ross2022molformer}). Table~\ref{tab:diff}
differentiates this work from the closest prior studies.

\begin{table}[t]\centering\scriptsize\setlength{\tabcolsep}{4pt}
\caption{Differentiation from the closest prior work. Y = yes, P = partial, -- = no.}
\label{tab:diff}
\begin{tabular}{lccccccc}
\toprule
Work & \#Endpts & Model fams & Conformal & Ensemble & Realistic split & Selective-pred & Found.\ model \\
\midrule
CFR \citep{li2024conformalized}      & P & 1 & Y & Y & P & P & -- \\
CoDrug \citep{laghuvarapu2023codrug} & P & 1 & Y & -- & Y & -- & -- \\
MOOD \citep{tossou2024mood}          & P & P & -- & Y & Y & -- & -- \\
Fooladi \citep{fooladi2025evaluating}& P & Y & -- & -- & Y & -- & -- \\
Koleiev \citep{koleiev2026critical}  & Y (22) & Y & -- & -- & P & -- & P \\
\textbf{This work}                   & \textbf{Y (22)} & \textbf{P} & \textbf{Y} & \textbf{Y} & \textbf{Y} & \textbf{Y} & \textbf{Y} \\
\bottomrule
\end{tabular}
\end{table}

\section{Data and Methods}
\textbf{Datasets.} All 22 endpoints of the TDC ADMET benchmark group (9 regression, 13
classification), via PyTDC 1.1.15 with pinned RDKit 2023.9.6. \textbf{Two loading modes.} For
leaderboard reproduction we use TDC's native, unpooled split and the official 5-seed
\texttt{admet\_group} evaluate protocol. For the controlled split-type comparison we pool
train+test, canonicalize, de-duplicate on canonical SMILES, and re-split under explicit seeds;
pooled numbers are internal-comparison-only. \textbf{Splits.} \emph{random} (optimistic baseline);
\emph{scaffold} (Bemis--Murcko, acyclic molecules treated as singleton groups, not lumped);
\emph{cluster} (sphere-exclusion via RDKit LeaderPicker on ECFP4---our primary OOD stressor,
$O(N\!\cdot\!L)$ and memory-safe on the $\sim$13k-molecule CYP sets). \textbf{Leakage} is reported
with a splitter-independent nearest-neighbour ECFP4 Tanimoto metric. \textbf{Models.} ECFP4 +
XGBoost baselines for the atlas, with a random-forest robustness spot-check on 5 endpoints; and a
frozen ChemBERTa-embedding arm with an XGBoost head (a representation swap, not a new model class). A
GNN baseline is left to future work.

\textbf{Generalization gap and statistics.} For each (dataset, model) and realistic split we compute
the gap per \emph{matched replicate} (same split seed and model seed for random and realistic, so the
comparison is paired). We report the mean gap, a 95\% bootstrap CI over replicates, a paired Wilcoxon
$p$-value, and Benjamini--Hochberg $q$-values across the 22-endpoint family; a gap whose CI crosses
zero is ``not established.'' An aggregate Wilcoxon over per-endpoint gaps gives the headline claim.

\textbf{Conformal and selective prediction.} Split-conformal with the validation fold as calibration:
absolute-residual intervals (regression) and LAC sets (classification). We implement weighted
conformal under covariate shift \citep{tibshirani2019conformal} with importance weights from an
isotonic-calibrated, class-balanced domain classifier, clipped at the 99th percentile, with an
effective-sample-size fallback. Our quantile reduces \emph{exactly} to textbook split-conformal under
uniform weights (unit-tested) and returns $+\infty$ when coverage cannot be certified. Selective
prediction uses risk--coverage curves with two label-free signals---deep-ensemble disagreement and
applicability-domain (AD) similarity---compared by Excess-AURC (AURC minus the oracle ordering),
which is comparable across folds with different base error rates.

\section{Results}

\subsection{The harness matches untuned fingerprint baselines}
Under the official 5-seed protocol, our deliberately simple ECFP4+XGBoost baselines land in the range
of TDC's untuned fingerprint baselines (Table~\ref{tab:lb}, full list in Appendix). They are
\emph{not} competitive with tuned leaderboard entries (e.g.\ our Caco-2 MAE 0.386 vs.\ a tuned XGBoost
leaderboard entry near 0.27); the point is only to confirm the harness is sane and in-range before any
re-split is introduced, so the subsequent re-split numbers are interpretable rather than an artifact.

\begin{table}[t]\centering\footnotesize
\caption{Leaderboard reproduction (official TDC split, 5-seed protocol), excerpt.}
\label{tab:lb}
% \begin{tabular}{llr}
\toprule
Endpoint & Metric & ECFP4+XGBoost (mean$\pm$std) \\
\midrule
ames & auroc & 0.832 $\pm$ 0.001 \\
bbb\_martins & auroc & 0.872 $\pm$ 0.006 \\
bioavailability\_ma & auroc & 0.693 $\pm$ 0.006 \\
caco2\_wang & mae & 0.386 $\pm$ 0.030 \\
clearance\_hepatocyte\_az & spearman & 0.316 $\pm$ 0.017 \\
clearance\_microsome\_az & spearman & 0.425 $\pm$ 0.021 \\
cyp2c9\_substrate\_carbonmangels & auprc & 0.371 $\pm$ 0.012 \\
cyp2c9\_veith & auprc & 0.753 $\pm$ 0.003 \\
cyp2d6\_substrate\_carbonmangels & auprc & 0.633 $\pm$ 0.016 \\
cyp2d6\_veith & auprc & 0.633 $\pm$ 0.004 \\
cyp3a4\_substrate\_carbonmangels & auroc & 0.641 $\pm$ 0.010 \\
cyp3a4\_veith & auprc & 0.852 $\pm$ 0.002 \\
dili & auroc & 0.857 $\pm$ 0.004 \\
half\_life\_obach & spearman & 0.301 $\pm$ 0.028 \\
herg & auroc & 0.727 $\pm$ 0.008 \\
hia\_hou & auroc & 0.955 $\pm$ 0.002 \\
ld50\_zhu & mae & 0.634 $\pm$ 0.011 \\
lipophilicity\_astrazeneca & mae & 0.666 $\pm$ 0.004 \\
pgp\_broccatelli & auroc & 0.908 $\pm$ 0.004 \\
ppbr\_az & mae & 10.270 $\pm$ 0.224 \\
solubility\_aqsoldb & mae & 1.254 $\pm$ 0.009 \\
vdss\_lombardo & spearman & 0.401 $\pm$ 0.016 \\
\bottomrule
\end{tabular}
  % full 22-row table auto-generated; excerpt shown
\begin{tabular}{llr}
\toprule Endpoint & Metric & ECFP4+XGBoost (mean$\pm$std) \\ \midrule
caco2\_wang & MAE & 0.386 $\pm$ 0.030 \\
bbb\_martins & AUROC & 0.872 $\pm$ 0.006 \\
herg & AUROC & 0.727 $\pm$ 0.008 \\
solubility\_aqsoldb & MAE & 1.254 $\pm$ 0.009 \\
\bottomrule
\end{tabular}
\end{table}

\subsection{Scaffold under-stresses; cluster reveals a universal gap}
Across all 22 endpoints (8 split seeds $\times$ 3 model seeds, matched-replicate \emph{relative} gaps
$=$ gap$/$random-split score; CIs from a split-seed block bootstrap that resamples the unit of
independence), under \emph{cluster} splits 21/22 endpoints show a statistically established gap (95\%
CI excludes zero), with median relative gap $+0.234$ and aggregate Wilcoxon $p<10^{-4}$. Under
\emph{scaffold} splits gaps are markedly smaller (median $+0.102$) and 5/22 are not established; most
strikingly hERG is flat under scaffold ($-0.002$, n.s.) but a clear $+0.066$ under cluster. The
cluster effect is $\approx$2.3$\times$ the scaffold median (Figure~\ref{fig:atlas}; per-endpoint
Appendix table with block-bootstrap CIs and per-split established flags). Because raw-unit gaps are
not cross-endpoint comparable (PPBR's 0--100\% target gives a large MAE gap), we report relative gaps
throughout. The realistic splits genuinely push test chemistry out of domain: mean nearest-neighbour
Tanimoto similarity to the training set falls from $0.56$ (random) to $0.41$ (scaffold) to $0.28$
(cluster) across the 22 endpoints (Figure~\ref{fig:leak}).

\begin{figure}[t]\centering
\includegraphics[width=0.55\linewidth]{fig5_leakage_gradient.png}
\caption{Mean nearest-neighbour Tanimoto similarity (test $\to$ nearest training molecule) per
endpoint. Realistic splits push the test set out of the training domain; cluster is the strongest
stressor.}
\label{fig:leak}
\end{figure}

\begin{figure}[t]\centering
\includegraphics[width=0.62\linewidth]{fig1_atlas.png}
\caption{Relative generalization gap per endpoint: cluster (red) exceeds scaffold (blue) on all 22
endpoints. hERG, flat under scaffold, shows a positive cluster gap.}
\label{fig:atlas}
\end{figure}

\subsection{The gap is an applicability-domain effect}
Binning every test prediction by nearest-neighbour Tanimoto similarity to training, normalized error
\emph{spikes} to $\approx$7.9$\times$ the endpoint mean in the most dissimilar bin (similarity $<0.1$,
a sparsely populated tail) and is already near-baseline ($\approx$1.3$\times$) by the next decile
(Figure~\ref{fig:ad}). Error is concentrated in molecules far from training---explaining the gap and
motivating distance as an abstention signal. Across endpoints, the drop in mean similarity from random
to cluster splits is suggestively associated with the relative gap (Spearman $\rho=+0.41$, $p=0.06$).
We flag a shared-representation caveat: ECFP4 Tanimoto defines the cluster split, the leakage metric,
\emph{and} this AD signal, so the relationship is partly built in; a splitter-independent distance
control is left to future work.

\begin{figure}[t]\centering
\includegraphics[width=0.55\linewidth]{fig2_applicability.png}
\caption{Normalized error vs.\ nearest-neighbour similarity to training (mean over 22 endpoints).}
\label{fig:ad}
\end{figure}

\subsection{Uncertainty under shift: an honest accounting}
Corrected split-conformal achieves nominal coverage on exchangeable data (self-test 0.901 at the 0.90
target). On real splits, coverage degradation under cluster shift is \emph{endpoint-dependent and
modest in aggregate}: mean cluster coverage is 0.896 (vs.\ 0.907 under random; target 0.90), but it
drops sharply on specific high-shift endpoints (BBB $0.92\!\rightarrow\!0.84$; 4/22 endpoints below
0.88; Figure~\ref{fig:cov}). Covariate-shift weighting is correctly a no-op on random splits but did \emph{not} restore
cluster coverage: the importance weights stayed near-uniform (Kish ESS $\approx$0.90), consistent with
the known weakness of classifier density-ratios under high-dimensional shift; we report ESS and fall
back rather than over-claim.

\begin{figure}[t]\centering
\includegraphics[width=0.55\linewidth]{fig6_coverage.png}
\caption{Split-conformal coverage per endpoint (target 0.90). Near-nominal on average for all splits,
but the lower tail degrades under cluster shift (worst: BBB $\approx$0.84); degradation is
endpoint-specific, not uniform.}
\label{fig:cov}
\end{figure}

For abstention, a random-forest pilot suggested AD distance beat ensemble disagreement under cluster
shift; \emph{this did not replicate} with XGBoost ensembles, where ensemble disagreement is
significantly better (wins on 18/22 endpoints, paired Wilcoxon $p=0.007$) and AD abstention actively
\emph{hurts} on 9/22. The headline is thus the cross-model \emph{instability} of the best abstention
signal, not a winner. Finally, we find no evidence that high-gap endpoints benefit more from abstention
(Spearman $\rho=+0.07$, 95\% CI $[-0.36,+0.48]$, under-powered at $n{=}22$; Figure~\ref{fig:unify}):
selective prediction does not cleanly ``recover the gap.''

\begin{figure}[t]\centering
\includegraphics[width=0.5\linewidth]{fig3_unifying.png}
\caption{Per-endpoint optimism gap vs.\ abstention benefit under cluster shift: no correlation
($\rho=0.07$, n.s.). An honest null.}
\label{fig:unify}
\end{figure}

\subsection{Do foundation models close the gap? Modestly}
Running the same machinery on frozen ChemBERTa embeddings + an XGBoost head, the foundation model
shrinks the relative cluster gap versus ECFP4 on 17/22 endpoints (mean $\Delta=-0.014$, paired
Wilcoxon $p=0.011$), but the effect is small and non-uniform, and the gap remains large
(Figure~\ref{fig:fm}; Table in Appendix). Pretrained representations help but do \emph{not} solve
out-of-distribution generalization.

\begin{figure}[t]\centering
\includegraphics[width=0.46\linewidth]{fig4_foundation.png}
\caption{ChemBERTa vs.\ ECFP4 relative cluster gap. Points below the diagonal = smaller gap.}
\label{fig:fm}
\end{figure}

\section{Discussion}
The message is not ``models are bad'' but ``the standard evaluation is optimistic in a
\emph{predictable} way'': measured performance tracks how far test chemistry sits from training, and
scaffold splits under-state it. The natural remedy---abstain on the OOD tail---is more fragile than
one would hope: conformal itself under-covers under the same shift, the best abstention signal is
model-dependent, and abstention benefit does not track the gap. The practical takeaway is to
\textbf{report cluster-split performance and nearest-neighbour applicability domain as standard}, and
to treat uncertainty methods as helpful but shift-sensitive rather than a turnkey fix. Our negative
results are prescriptive: classifier density-ratio covariate-shift correction is unreliable for ADMET
fingerprints (ESS collapse) and should not be relied upon, and the best abstention signal is
model-dependent---there is no turnkey uncertainty signal, so it must be validated per model class.

\section{Limitations}
Compute was bounded to a single consumer device (Apple M4 Pro), which keeps the study fully
reproducible on commodity hardware but limits the foundation-model arm to frozen embeddings; GNNs and
fine-tuning are future work. Classifier density-ratio weights collapse in ESS under severe shift (we
report ESS and fall back). Temporal splits are out of scope for most endpoints (no provenance dates)
and not claimed. Pooled re-splits are not leaderboard-comparable by construction; we anchor with the
official-split reproduction.

\section{Conclusion}
A reproducible, leakage-audited re-evaluation of TDC ADMET shows \emph{when} and \emph{why} the
generalization gap appears---universally under cluster splits, driven by distance-to-training---and
gives an honest accounting of what uncertainty methods do and do not buy under shift. Code, pinned
environment, splits, and figures are released.

\small
\bibliographystyle{plainnat}
\bibliography{references}

\normalsize
\appendix
\section{Full leaderboard reproduction (all 22 endpoints)}
\begin{center}\footnotesize\begin{tabular}{llr}
\toprule
Endpoint & Metric & ECFP4+XGBoost (mean$\pm$std) \\
\midrule
ames & auroc & 0.832 $\pm$ 0.001 \\
bbb\_martins & auroc & 0.872 $\pm$ 0.006 \\
bioavailability\_ma & auroc & 0.693 $\pm$ 0.006 \\
caco2\_wang & mae & 0.386 $\pm$ 0.030 \\
clearance\_hepatocyte\_az & spearman & 0.316 $\pm$ 0.017 \\
clearance\_microsome\_az & spearman & 0.425 $\pm$ 0.021 \\
cyp2c9\_substrate\_carbonmangels & auprc & 0.371 $\pm$ 0.012 \\
cyp2c9\_veith & auprc & 0.753 $\pm$ 0.003 \\
cyp2d6\_substrate\_carbonmangels & auprc & 0.633 $\pm$ 0.016 \\
cyp2d6\_veith & auprc & 0.633 $\pm$ 0.004 \\
cyp3a4\_substrate\_carbonmangels & auroc & 0.641 $\pm$ 0.010 \\
cyp3a4\_veith & auprc & 0.852 $\pm$ 0.002 \\
dili & auroc & 0.857 $\pm$ 0.004 \\
half\_life\_obach & spearman & 0.301 $\pm$ 0.028 \\
herg & auroc & 0.727 $\pm$ 0.008 \\
hia\_hou & auroc & 0.955 $\pm$ 0.002 \\
ld50\_zhu & mae & 0.634 $\pm$ 0.011 \\
lipophilicity\_astrazeneca & mae & 0.666 $\pm$ 0.004 \\
pgp\_broccatelli & auroc & 0.908 $\pm$ 0.004 \\
ppbr\_az & mae & 10.270 $\pm$ 0.224 \\
solubility\_aqsoldb & mae & 1.254 $\pm$ 0.009 \\
vdss\_lombardo & spearman & 0.401 $\pm$ 0.016 \\
\bottomrule
\end{tabular}
\end{center}

\section{Per-endpoint generalization-gap atlas}
Relative gaps (gap $/$ random-split score); cluster column with 95\% bootstrap CI and BH-FDR $q$.
\begin{center}\footnotesize\begin{tabular}{llrrl}
\toprule
Endpoint & Metric & Scaffold & Cluster & Cluster 95\% CI ($q$) \\
\midrule
ppbr\_az & mae & 0.229 & 0.583 & [+5.614, +5.909] ($q$0.000) \\
clearance\_microsome\_az & spearman & 0.314 & 0.548 & [+0.261, +0.304] ($q$0.000) \\
lipophilicity\_astrazeneca & mae & 0.216 & 0.525 & [+0.307, +0.318] ($q$0.000) \\
clearance\_hepatocyte\_az & spearman & 0.195 & 0.519 & [+0.222, +0.273] ($q$0.000) \\
vdss\_lombardo & spearman & 0.114 & 0.483 & [+0.219, +0.258] ($q$0.000) \\
half\_life\_obach & spearman & 0.232 & 0.480 & [+0.106, +0.186] ($q$0.000) \\
solubility\_aqsoldb & mae & 0.216 & 0.425 & [+0.404, +0.425] ($q$0.000) \\
caco2\_wang & mae & 0.287 & 0.311 & [+0.092, +0.111] ($q$0.000) \\
cyp2d6\_veith & auprc & 0.103 & 0.305 & [+0.204, +0.219] ($q$0.000) \\
cyp2d6\_substrate\_carbonmangels & auprc & 0.143 & 0.247 & [+0.107, +0.196] ($q$0.000) \\
ld50\_zhu & mae & 0.101 & 0.244 & [+0.108, +0.118] ($q$0.000) \\
cyp2c9\_substrate\_carbonmangels & auprc & 0.147 & 0.225 & [+0.058, +0.102] ($q$0.000) \\
cyp3a4\_veith & auprc & 0.058 & 0.225 & [+0.188, +0.199] ($q$0.000) \\
bbb\_martins & auroc & 0.020 & 0.186 & [+0.164, +0.177] ($q$0.000) \\
pgp\_broccatelli & auroc & 0.029 & 0.185 & [+0.166, +0.177] ($q$0.000) \\
ames & auroc & 0.072 & 0.160 & [+0.137, +0.147] ($q$0.000) \\
cyp2c9\_veith & auprc & 0.049 & 0.159 & [+0.119, +0.133] ($q$0.000) \\
dili & auroc & 0.030 & 0.106 & [+0.069, +0.111] ($q$0.000) \\
bioavailability\_ma & auroc & 0.055 & 0.076 & [+0.035, +0.073] ($q$0.000) \\
cyp3a4\_substrate\_carbonmangels & auroc & 0.034 & 0.073 & [+0.030, +0.066] ($q$0.000) \\
herg & auroc & -0.002 & 0.066 & [+0.043, +0.068] ($q$0.000) \\
hia\_hou & auroc & 0.022 & 0.045 & [+0.016, +0.070] ($q$0.010) \\
\bottomrule
\end{tabular}
\end{center}

\section{Foundation model vs.\ fingerprint (relative cluster gap)}
\begin{center}\footnotesize\begin{tabular}{lrrr}
\toprule
Endpoint & ECFP4 & ChemBERTa & $\Delta$ \\
\midrule
ppbr\_az & 0.583 & 0.443 & -0.140 \\
vdss\_lombardo & 0.483 & 0.357 & -0.125 \\
lipophilicity\_astrazeneca & 0.525 & 0.419 & -0.106 \\
clearance\_microsome\_az & 0.548 & 0.473 & -0.075 \\
cyp2d6\_substrate\_carbonmangels & 0.247 & 0.184 & -0.064 \\
ld50\_zhu & 0.244 & 0.183 & -0.061 \\
cyp2c9\_veith & 0.159 & 0.112 & -0.047 \\
pgp\_broccatelli & 0.185 & 0.142 & -0.044 \\
bioavailability\_ma & 0.076 & 0.038 & -0.038 \\
bbb\_martins & 0.186 & 0.155 & -0.031 \\
dili & 0.106 & 0.076 & -0.030 \\
cyp3a4\_substrate\_carbonmangels & 0.073 & 0.052 & -0.021 \\
clearance\_hepatocyte\_az & 0.519 & 0.503 & -0.016 \\
solubility\_aqsoldb & 0.425 & 0.410 & -0.015 \\
cyp2d6\_veith & 0.305 & 0.293 & -0.013 \\
cyp3a4\_veith & 0.225 & 0.215 & -0.010 \\
herg & 0.066 & 0.060 & -0.006 \\
hia\_hou & 0.045 & 0.045 & +0.000 \\
ames & 0.160 & 0.160 & +0.001 \\
caco2\_wang & 0.311 & 0.320 & +0.009 \\
half\_life\_obach & 0.480 & 0.721 & +0.241 \\
cyp2c9\_substrate\_carbonmangels & 0.225 & 0.502 & +0.277 \\
\midrule
Mean & 0.281 & 0.267 & -0.014 \\
\bottomrule
\end{tabular}
\end{center}
